\section{詳細解説}

\subsection{数値例題：二本脚の膝トルク}

\begin{example}[20 cm の段差]
股関節の高さ $H=0.9$ m、大腿・下腿 $\ell=0.45$ m、段差 $h=0.2$ m とすると、段の上の足首から股関節までの距離は $H-h=0.7$ m。
膝のはみ出しは $d=\sqrt{0.45^2-0.35^2} = 0.283$ m。膝より上の質量 $m=93$ kg を、上体を立てたまま（$x_c=0$）持ち上げるには
\[
\tau_{\text{膝}} = 93\times9.81\times0.283 = 258\ \text{N\,m}
\]
が必要で、上限 250 N\,m をわずかに超える（第~\ref{ch:steps}章の表の値と一致する）。
上体を前に傾けて重心を足首から 0.15 m 前に移すと、$\tau_{\text{膝}} = 93\times9.81\times(0.283-0.15) = 121$ N\,m、
$\tau_{\text{股}} \approx 93\times9.81\times0.15 = 137$ N\,m となり、どちらも上限内に収まる。
\end{example}

\subsection{感度分析}

段差の限界がどのパラメータで決まるかを調べた（図~\ref{fig:D13_step_sensitivity.png}）。
\begin{itemize}
\item \textbf{戦車型と摩擦係数}：力の限界（腕を前に伸ばした場合）は、摩擦係数 0.2 以上では 0.33 m で一定だった。
  この範囲では、限界を決めているのは摩擦ではなく、後ろへの転倒（重心が後ろの接地点を越える）である。
  摩擦係数が 0.15 以下（氷や濡れた路面）では、限界が 0.08 m 以下に急に下がり、摩擦が支配的になる。
  前端の高さ（0.25 m）の形の限界は、この 0.33 m より小さいので、通常の路面では形が限界を決める。
\item \textbf{二本脚と関節トルク}：限界の段差は、膝・股関節のトルク上限にほぼ比例して増え、250〜275 N\,m の間で可動域の限界（0.63 m）を超える。
  人間大の二本脚では、\textbf{トルクと可動域の両方がちょうど釣り合う設計点}にあると言える。
\end{itemize}

\DFIG{D13_step_sensitivity.png}{段差の限界の感度。左：戦車型の力の限界と摩擦係数（破線は前端の高さ）。右：二本脚の限界と膝・股関節のトルク上限（破線は可動域の限界）。}

\subsection{相似拡大の全結果}

\DTAB{sim10_scaling}

\subsection{なぜ戦車型はスケール不変なのか}

戦車型の段差の限界は、(1) 前端の高さ（形）、(2) 2 点支持のつり合いと摩擦（角度と力の比）で決まる。
どちらも長さの比と無次元の摩擦係数だけで決まり、重さや大きさの絶対値には依存しない。
したがって、相似拡大すると越えられる段差も同じ比率で大きくなる（大きさの 1/4）。
一方、二本脚は「体を持ち上げる力（トルク）」で決まるため、二乗三乗則の影響をまともに受ける。

\section*{考察課題}
\begin{exercise}
二本脚が勢いを使って（蹴り出して）段差を上る場合、必要なのはトルクではなく何か。大型化したときの不利は残るか。
\end{exercise}
\begin{exercise}
戦車型にフリッパー（前端を持ち上げる補助の履帯）を付けて前端の高さを 0.5 m にすると、越えられる段差は何で決まるようになるか。図~\ref{fig:D13_step_sensitivity.png} から考えよ。
\end{exercise}
